% ============================================================
% CONCLUSION AND FUTURE WORK
% ============================================================
\section{Conclusion and Future Work}
\label{sec:conclusion}

Current language models did not reproduce the racial or sex disparities
documented in human pain assessment. Four open-weight models showed none on six
authored cases, whether race and sex were stated or implied by a name, and the
one apparent pattern was an artefact of a single name. A panel of
\nConfigs{} proprietary and open-weight configurations showed none on 50
clinician-written cases either: severity ratings were equivalent across
identities within one point, no difference survived the study-wide false
discovery rate, and the small differences that survived within a model favoured
Black patients. What the models were not was interchangeable. They disagreed
with one another and with the clinical reference on whether to offer an opioid
and at what dose, and the same weights answered differently with reasoning on or
off and at a different precision, by margins larger than any identity effect.
Beyond the substantive null, the study contributes an audit protocol for
clinical LLMs: paired seeds, several names per group, the case as the unit of
analysis, a reported noise floor and minimum detectable effect, and a separate
audit for each deployed configuration.

Three extensions follow directly. The cross-lingual audit, in the four PAINT
languages \cite{zhang2026paint}, tests whether alignment that holds in English
holds elsewhere; the request code and draft translations are already in the
repository. Richer cues, including dialect \cite{hofmann2024dialect},
stigmatising language \cite{huang2026artificial}, insurance status and
realistic note length, test whether the null survives the way identity actually
appears in records. Free-text responses, realistic record context and
outcome-grounded designs \cite{jarrett2026outcome} test whether the constraint
of a short JSON answer is what keeps the disparity out, and whether the
between-model differences in dosing translate into differences in outcome.
